% Requer no preâmbulo: \usepackage{booktabs}, \usepackage{graphicx} e \usepackage{tabularx}
\begin{table*}[htbp]
  \centering
  \caption{Resultados comparativos do \texttt{rtpjitterbuffer} nos seis cenários experimentais (Desktop Ethernet \texttt{enp37s0} $\rightarrow$ Notebook Wi-Fi).}
  \label{tab:resultados_cenarios}
  \resizebox{\textwidth}{!}{%
  \begin{tabular}{llrrrrrrrrl}
    \toprule
    \textbf{Cen.} & \textbf{Rede (\texttt{tc netem})} & \textbf{Lat. (ms)} & \textbf{Dur. (s)} & \textbf{Receb.} & \textbf{Perd.} & \textbf{Atras.} & \textbf{Dup.} & \textbf{Perda (\%)} & \textbf{Jitter médio $\pm$ DP (ms)} & \textbf{Mediana [mín--máx] (ms)} \\
    \midrule
    C1 & Normal (sem netem) & 50 & 64 & 15 270 & 2 & 0 & 0 & 0,01\% & 9,06 $\pm$ 1,91 & 8,25 [7,86--15,65] \\
    C2 & Normal (sem netem) & 200 & 66 & 15 839 & 0 & 0 & 0 & 0,00\% & 8,62 $\pm$ 1,13 & 8,23 [8,08--12,84] \\
    C3 & delay 50ms 40ms & 50 & 82 & 18 385 & 1 098 & 1 089 & 9 & 5,64\%* & 25,25 $\pm$ 8,69 (29,34 $\pm$ 2,43*) & 28,59 [8,26--35,26] \\
    C4 & delay 50ms 40ms & 200 & 67 & 15 821 & 5 & 3 & 0 & 0,03\% & 29,60 $\pm$ 2,92 & 29,70 [24,10--35,15] \\
    C5 & loss 20\% delay 50ms 40ms & 50 & 74 & 12 283 & 5 063 & 1 549 & 5 & 29,19\% & 29,46 $\pm$ 3,55 & 29,33 [22,75--38,66] \\
    C6 & loss 20\% delay 50ms 40ms & 1000 & 66 & 12 401 & 3 125 & 0 & 0 & 20,13\% & 29,27 $\pm$ 3,52 & 29,32 [20,30--35,43] \\
    \bottomrule
  \end{tabular}}
  \vspace{4pt}
  \parbox{\textwidth}{\footnotesize \textit{*Em C3, a regra tc netem foi ativada após o arranque do vídeo (aos 18 s, passando o jitter de 8,44 ms para 29,25 ms). Isolando apenas a janela de regime estacionário (18 s a 82 s), o jitter médio foi de 29,34 $\pm$ 2,43 ms e a taxa de perda efetiva foi de 6,84\% (1 045 perdas, das quais 1 037 --- ou 99,23\% --- por atraso).}}
\end{table*}

\begin{table*}[htbp]
  \centering
  \caption{Síntese dos principais achados experimentais e decomposição de perdas (C1--C6).}
  \label{tab:principais_achados}
  \begin{tabularx}{\textwidth}{lccccX}
    \toprule
    \textbf{Dimensão} & \textbf{Cenários} & \textbf{Perda Real} & \textbf{Atraso (\texttt{late})} & \textbf{Perda Total} & \textbf{Principal Achado} \\
    \midrule
    1. Linha de base Ethernet $\rightarrow$ Wi-Fi & C1 vs. C2 (50 vs. 200 ms) & 0,01\% $\rightarrow$ 0,00\% & 0 $\rightarrow$ 0 (0,00\%) & 0,01\% $\rightarrow$ 0,00\% & Jitter estável (8,62--9,06 ms) e integridade quase total sem degradação artificial (late = 0). \\
    2. Absorção de jitter puro ($\pm$40 ms) & C3 vs. C4 (50 vs. 200 ms) & 0,05\% $\rightarrow$ 0,01\% & 1 089 $\rightarrow$ 3 (-99,72\%) & 5,64\%* $\rightarrow$ 0,03\% (-99,44\%) & Em 50 ms, 99,2\% das perdas foram por atraso; o buffer de 200 ms absorveu o jitter mantendo baixa latência. \\
    3. Perda 20\% + jitter ($\pm$40 ms) & C5 vs. C6 (50 vs. 1000 ms) & 20,26\% $\rightarrow$ 20,13\% & 1 549 (8,93\%) $\rightarrow$ 0 (0,00\%) & 29,19\% $\rightarrow$ 20,13\% (-9,06 p.p.) & Buffer de 1 000 ms zerou descartes por atraso (late = 0), convergindo ao piso físico de 20\% do RTP/UDP sem retransmissão. \\
    \bottomrule
  \end{tabularx}
\end{table*}
